\chapter{Contexte et revue de littérature}
\label{chapitre:contexte}

% =====================================================================
\section{Cadre général}
% =====================================================================

La boiterie bovine est un problème de santé majeur en élevage laitier, avec des conséquences
directes sur le bien-être animal, la productivité et la rentabilité des exploitations.
Des études épidémiologiques rapportent des prévalences allant de 20\% à plus de 50\% selon
les régions, les systèmes d'élevage et les critères de diagnostic utilisés
\citep{whay2003assessment,cook2003prevalence}. Les pathologies du pied, qu'elles soient
infectieuses (\textit{dermatite digitale}, \textit{phlegmon interdigité}) ou non
infectieuses (\textit{ulcère de la sole}, \textit{maladie de la ligne blanche}),
constituent la cause la plus fréquente de boiterie chez les vaches laitières
\citep{merck_hoof_origin_lameness_cattle,flower2006effect}.

\subsection{Impact économique}
L'impact économique de la boiterie est substantiel et multidimensionnel. Le
Tableau~\ref{tab:impact_economique} synthétise les principales composantes du coût
\textit{par cas clinique}, c'est-à-dire par épisode de boiterie diagnostiqué chez une
vache au cours d'une lactation, d'après les estimations publiées dans la littérature.

\begin{table}[H]
    \centering
    \caption{Composantes du coût économique de la boiterie par cas clinique.}
    \label{tab:impact_economique}
    \begin{tabular}{lcc}
        \toprule
        \textbf{Composante} & \textbf{Estimation} & \textbf{Source} \\
        \midrule
        Perte de production laitière & 270--574 kg/lactation & \cite{green2002financial} \\
        Coût vétérinaire direct & 25--90 USD/cas & \cite{bruijnis2010assessing} \\
        Dégradation de la fertilité & +30--50 jours ouverts & \cite{green2002financial} \\
        Réforme anticipée & 75--300 USD/cas & \cite{bruijnis2010assessing} \\
        \midrule
        \textbf{Coût total estimé} & \textbf{75--300+ USD/cas} & \\
        \bottomrule
    \end{tabular}
\end{table}
\vspace{-0.9\baselineskip}

Quatre postes principaux composent le coût d'un cas clinique: la perte de production
laitière (270 à 574~kg par lactation), liée à la baisse de l'ingestion et de l'activité;
le coût vétérinaire direct (25 à 90~USD par cas), qui inclut la consultation, le parage,
les traitements et le matériel de soin; la dégradation de la fertilité, mesurée par une
hausse de 30 à 50 \textit{jours ouverts}, soit l'intervalle entre le vêlage et la
conception suivante, avec un décalage de la lactation suivante; et la réforme anticipée
(75 à 300~USD par cas), liée à la perte de valeur résiduelle d'une vache écartée avant la
fin de sa carrière productive. Ces postes se cumulent selon la sévérité et la chronicité
de l'épisode. À l'échelle d'un troupeau de 100~vaches avec une prévalence de 25\,\%, le
coût annuel cumulé est estimé entre 1\,900 et 7\,500~USD pour les seuls cas incidents.
Au-delà de l'impact financier, la boiterie soulève un enjeu éthique de bien-être animal~:
la douleur associée est reconnue comme l'une des principales sources de souffrance
chronique chez les vaches laitières \citep{whay2003assessment}. Ces conséquences
justifient une compréhension fine des mécanismes de boiterie, afin d'orienter la
prévention, la détection et la prise en charge.

\subsection{Étiologie et facteurs de risque}
Les causes de la boiterie sont multifactorielles. La Figure~\ref{fig:etiologie} les
organise en quatre catégories~: causes directes infectieuses ou non infectieuses, et
facteurs modulateurs environnementaux ou individuels. Les causes directes regroupent des
affections bactériennes, comme la \textit{dermatite digitale} ou le \textit{phlegmon
interdigité}, et des lésions mécaniques, comme l'\textit{ulcère de la sole} ou la
\textit{maladie de la ligne blanche}. Les facteurs environnementaux (sol, hygiène,
humidité) modulent l'exposition, tandis que les facteurs individuels (parité, âge, stade
de lactation) modulent la susceptibilité. Cette structure multifactorielle explique
qu'aucun indicateur unique ne suffise et justifie une approche multi-signaux pour la
détection précoce.

\begin{figure}[H]
    \centering
    \resizebox{0.85\textwidth}{!}{%
    \begin{tikzpicture}[
        cat/.style={rectangle, draw=black!70, rounded corners=3pt, minimum height=0.9cm,
                    minimum width=3.2cm, align=center, font=\small, inner sep=10pt},
        factor/.style={rectangle, draw=black!40, rounded corners=2pt, minimum height=0.8cm,
                      minimum width=2.7cm, align=center, font=\small, inner sep=6pt},
        arrow/.style={-{Stealth[length=2mm]}, thick, black!60},
    ]
        % Titre central
        \node[cat, fill=rawred!15, minimum width=4cm, font=\small\bfseries] (boiterie)
            {Boiterie bovine};

        % Catégories
        \node[cat, fill=ifblue!12, above left=1.5cm and 1cm of boiterie] (infect)
            {Pathologies\\[-7pt]infectieuses};
        \node[cat, fill=loforange!12, above right=1.5cm and 1cm of boiterie] (noninfect)
            {Pathologies\\[-7pt]non infectieuses};
        \node[cat, fill=rulegreen!12, below left=1.5cm and 1cm of boiterie] (enviro)
            {Facteurs\\[-7pt]environnementaux};
        \node[cat, fill=profgray!12, below right=1.5cm and 1cm of boiterie] (indiv)
            {Facteurs\\[-7pt]individuels};

        % Liens
        \draw[arrow] (infect) -- (boiterie);
        \draw[arrow] (noninfect) -- (boiterie);
        \draw[arrow] (enviro) -- (boiterie);
        \draw[arrow] (indiv) -- (boiterie);

        % Sous-facteurs infectieux
        \node[factor, fill=ifblue!6, above=0.9cm of infect, xshift=-2.0cm] (dd)
            {Dermatite digitale};
        \node[factor, fill=ifblue!6, above=0.9cm of infect, xshift=2.0cm] (phleg)
            {Phlegmon interdigité};
        \draw[black!30] (dd) -- (infect);
        \draw[black!30] (phleg) -- (infect);

        % Sous-facteurs non infectieux
        \node[factor, fill=loforange!6, above=0.9cm of noninfect, xshift=-2.0cm] (ulcere)
            {Ulcère de la sole};
        \node[factor, fill=loforange!6, above=0.9cm of noninfect, xshift=2.0cm] (ligne)
            {Maladie ligne blanche};
        \draw[black!30] (ulcere) -- (noninfect);
        \draw[black!30] (ligne) -- (noninfect);

        % Sous-facteurs environnementaux
        \node[factor, fill=rulegreen!6, below=0.9cm of enviro, xshift=-2.0cm] (sol)
            {Type de sol};
        \node[factor, fill=rulegreen!6, below=0.9cm of enviro, xshift=2.0cm] (hygiene)
            {Hygiène, humidité};
        \draw[black!30] (sol) -- (enviro);
        \draw[black!30] (hygiene) -- (enviro);

        % Sous-facteurs individuels
        \node[factor, fill=profgray!6, below=0.9cm of indiv, xshift=-2.0cm] (parite)
            {Parité, âge};
        \node[factor, fill=profgray!6, below=0.9cm of indiv, xshift=2.0cm] (lactation)
            {Stade de lactation};
        \draw[black!30] (parite) -- (indiv);
        \draw[black!30] (lactation) -- (indiv);
    \end{tikzpicture}%
    }
    \caption{Étiologie multifactorielle de la boiterie bovine.}
    \label{fig:etiologie}
\end{figure}

\subsection{Importance de la détection précoce}
La littérature sur la boiterie bovine souligne l'importance d'une détection précoce pour
limiter la dégradation du bien-être animal et les impacts de production. Les ressources
vétérinaires de référence et les guides techniques convergent sur deux idées:
une surveillance régulière est indispensable, et la prise en charge doit être
structurée et rapide selon la sévérité du cas
\citep{merck_lameness_overview,dairynz_treating_lameness}.

La détection précoce (score locomoteur 2 sur une échelle de 1 à 5) ouvre la voie au
traitement des cas avant qu'ils ne deviennent chroniques. Les études montrent qu'un
retard de traitement de quelques jours peut multiplier la durée de récupération et le
coût associé \citep{green2002financial}. Des travaux récents utilisant des capteurs
automatisés portés (podomètres et accéléromètres fixés à la patte ou au collier)
confirment que les systèmes d'alerte précoce peuvent identifier des changements
comportementaux deux à quatre jours avant la détection visuelle par l'éleveur
\citep{chapinal2010measurement,pastell2009measurement}. En pratique, les éleveurs
détectent souvent la boiterie à un stade avancé (score 3 ou plus), lorsque les signes
sont visibles à l'œil nu mais que le pronostic est déjà moins favorable.

Des revues de synthèse en élevage de précision montrent aussi que la boiterie fait partie
des cas d'usage les plus importants, mais également des plus difficiles à stabiliser, car
les signaux disponibles restent indirects, multifactoriels et fortement dépendants du
contexte d'élevage \citep{palma2025scoping,michelena2025plf_review}.

% =====================================================================
\section{Approches conventionnelles~: observation et notation locomotrice}
% =====================================================================

Les approches conventionnelles reposent principalement sur l'observation humaine et les
grilles de notation locomotrice. La grille proposée par \citet{sprecher1997lameness}
reste l'une des plus utilisées en pratique. Elle attribue un score de 1
(\textit{normal}) à 5 (\textit{sévèrement boiteux}) en se basant sur la posture du dos
et la démarche \citep{flower2006gait}. La fiche opérationnelle de l'AHDB propose une
grille complémentaire adaptée au contexte britannique et associe le score~2 à une
intervention rapide, généralement dans les 48~heures
\citep{ahdb_mobility_score_sheet_48h_score2}.

Les approches conventionnelles restent pertinentes en élevage car elles sont simples à
déployer et directement interprétables par les équipes terrain. Elles présentent toutefois
plusieurs limites opérationnelles. La première tient à la subjectivité inter-observateurs: même avec
des grilles standardisées, la concordance entre évaluateurs peut varier significativement
\citep{flower2006effect,schlageter2014validation}. La deuxième concerne le temps requis:
dans les grands troupeaux, l'observation systématique de chaque animal à chaque traite est
difficile à maintenir. Enfin, les boiteries légères (score~2) sont souvent sous-détectées
par l'observation visuelle, alors qu'elles représentent le stade où l'intervention est la
plus efficace.

Pour dépasser ces limites, le recours à des signaux capteurs vient compléter la
surveillance humaine, sans la remplacer.

\subsection{Lecture opérationnelle du score locomoteur}
Au-delà de sa fonction descriptive, l'évaluation locomotrice est aussi un outil de
priorisation de l'intervention. En pratique, la zone la plus critique est le passage du score~1
(\textit{normal}) au score~2 (\textit{boiterie légère}): c'est à ce stade que les signes
deviennent suffisamment perceptibles pour justifier une action, tout en restant assez précoces
pour espérer limiter l'installation d'une boiterie chronique. Cette logique explique
pourquoi de nombreux systèmes automatisés ne cherchent pas directement à reproduire
l'intégralité d'une note clinique, mais plutôt à signaler une \textit{boiterie probable}
suffisamment tôt pour déclencher une vérification sur le terrain.

La Figure~\ref{fig:sprecher_scale} synthétise la lecture opérationnelle du score
locomoteur: le passage vers le score~2 constitue la zone d'intérêt pour une alerte
précoce, car il marque l'entrée dans une situation où l'observation humaine ou l'examen
du pied devient prioritaire.

\begin{figure}[H]
    \centering
    \begin{tikzpicture}[
        score/.style={rectangle, draw=black!60, rounded corners=3pt, minimum width=2.45cm,
                      minimum height=1.7cm, align=center, inner sep=8pt, font=\small},
        arrow/.style={-{Stealth[length=2.5mm]}, thick, black!55}
    ]
        \node[score, fill=rulegreen!12] (s1) at (0,0.2) {\textbf{Score 1}\\[2pt]Normal};
        \node[score, fill=ifblue!10] (s2) at (3.2,0.2) {\textbf{Score 2}\\[2pt]Léger};
        \node[score, fill=loforange!14] (s3) at (6.4,0.2) {\textbf{Score 3}\\[2pt]Modéré};
        \node[score, fill=rawred!14] (s4) at (9.6,0.2) {\textbf{Score 4}\\[2pt]Marqué};
        \node[score, fill=rawred!25] (s5) at (12.8,0.2) {\textbf{Score 5}\\[2pt]Sévère};

        \draw[arrow] (1.35,0.2) -- (1.95,0.2);
        \draw[arrow] (4.55,0.2) -- (5.15,0.2);
        \draw[arrow] (7.75,0.2) -- (8.35,0.2);
        \draw[arrow] (10.95,0.2) -- (11.55,0.2);

        \draw[decorate, decoration={brace, amplitude=6pt}, thick, ifblue]
            (1.7,1.7) -- (4.7,1.7)
            node[midway, above=8pt, align=center, font=\footnotesize\bfseries]
            {Zone d'intérêt pour\\la détection précoce};
    \end{tikzpicture}

    \caption[Seuil opérationnel du score locomoteur]{Seuil opérationnel du score
    locomoteur pour la détection précoce.}
    \label{fig:sprecher_scale}
\end{figure}

Le Tableau~\ref{tab:score_operationnel} précise la lecture clinique simplifiée et
l'action associée à chaque score, d'après la grille de \citet{sprecher1997lameness} et
les recommandations AHDB \citep{ahdb_mobility_scoring_dairy_cows}.

\begin{table}[H]
    \centering
    \caption{Lecture clinique et opérationnelle simplifiée du score locomoteur.}
    \label{tab:score_operationnel}
    \footnotesize
    \renewcommand{\arraystretch}{1.2}
    \begin{tabular}{>{\raggedright\arraybackslash}p{1.4cm}
                    >{\raggedright\arraybackslash}p{6.0cm}
                    >{\raggedright\arraybackslash}p{5.2cm}}
        \toprule
        \textbf{Score} & \textbf{Lecture clinique simplifiée} & \textbf{Lecture opérationnelle} \\
        \midrule
        1 & Dos plat en station et en marche. & Surveillance usuelle. \\
        2 & Dos plat à l'arrêt, légèrement arqué en marche. & Intervention rapide: zone clé pour la détection précoce. \\
        3 & Dos arqué en permanence. & Vérification prioritaire. \\
        4 & Atteinte locomotrice visible. & Prise en charge urgente. \\
        5 & Appui très limité. & Prise en charge immédiate. \\
        \bottomrule
    \end{tabular}
\end{table}

La lecture par seuil transforme la notation locomotrice en problème d'alerte: il ne
s'agit plus seulement de décrire une boiterie, mais de repérer l'entrée dans une zone
d'intervention. Ce déplacement justifie le recours à des signaux capteurs et à des
méthodes d'apprentissage capables de suivre les changements comportementaux dans le
temps.

% =====================================================================
\section{Approches basées sur des capteurs et apprentissage}
% =====================================================================

Les approches basées sur des capteurs visent à détecter des déviations comportementales
(activité, transitions, dynamique de mouvement) compatibles avec une boiterie probable.
La littérature distingue trois familles principales: l'analyse de la démarche par vision
par ordinateur \citep{vanhertem2014evaluation,lavrova2023lameness}, la classification de
comportements par accélérométrie \citep{oleary2020accelerometers,alsaaod2019automatic,
balasso2023behaviour} et la détection d'anomalies sur des indicateurs dérivés
\citep{qiao2021review,gonzalezgarcia2023systematic,chandola2009anomaly}. Les revues
récentes confirment la prédominance des approches supervisées, tout en soulignant le
manque d'études non supervisées en conditions réelles d'élevage
\citep{qiao2021review,gonzalezgarcia2023systematic}.

La diversité de ces approches peut être organisée selon trois dimensions complémentaires:
\textit{la source de données}, \textit{la stratégie analytique} et \textit{la nature de la
sortie}. Une telle lecture est importante, car elle évite de comparer directement des travaux
qui ne visent pas le même objectif. Un modèle supervisé de classification binaire à partir
d'images et un système d'alerte basé sur des capteurs d'activité n'ont ni les mêmes
contraintes de données, ni le même niveau d'interprétabilité, ni le même usage en ferme.

La Figure~\ref{fig:taxonomie_approches} synthétise cette organisation en trois axes~: la
source de données (de l'observation visuelle aux accéléromètres portés), la stratégie
analytique (règles et seuils, modèles supervisés, détection d'anomalies non supervisée,
approches hybrides) et la sortie opérationnelle (du score locomoteur à l'alerte
actionnable). Le présent travail se positionne dans la branche non supervisée et hybride,
encadrée en vert sur la figure.

\begin{figure}[H]
    \centering
    \resizebox{\textwidth}{!}{%
    \begin{tikzpicture}[
        block/.style={rectangle, draw=black!70, rounded corners=3pt, minimum width=4cm,
                      minimum height=1cm, align=center, font=\small\bfseries, inner sep=8pt},
        small/.style={rectangle, draw=black!35, rounded corners=2pt, minimum width=2.8cm,
                      minimum height=0.85cm, text width=2.6cm, align=center, font=\footnotesize, inner sep=6pt},
        arrow/.style={-{Stealth[length=2.5mm]}, thick, black!55}
    ]
        \node[block, fill=ifblue!12] (data) at (0,0) {Sources de données};
        \node[block, fill=loforange!12] (model) at (7,0) {Stratégie analytique};
        \node[block, fill=rulegreen!12] (output) at (14,0) {Sortie opérationnelle};

        \node[small, fill=ifblue!6] (obs) at (-1.8,-2.2) {Observation\\et notation};
        \node[small, fill=ifblue!6] (vision) at (1.8,-2.2) {Vision\\de la démarche};
        \node[small, fill=ifblue!6] (accel) at (-1.8,-4.0) {Accéléromètres\\activité / couchage};
        \node[small, fill=ifblue!6] (multi) at (1.8,-4.0) {Multi-source\\capteurs + contexte};

        \node[small, fill=loforange!6] (rules) at (5.2,-2.2) {Règles\\et seuils};
        \node[small, fill=loforange!6] (sup) at (8.8,-2.2) {Supervisé};
        \node[small, fill=loforange!6] (unsup) at (5.2,-4.0) {Non supervisé\\(anomalies)};
        \node[small, fill=loforange!6] (hyb) at (8.8,-4.0) {Hybride\\score + règles};

        \node[small, fill=rulegreen!6] (score) at (12.2,-2.2) {Score\\locomoteur};
        \node[small, fill=rulegreen!6] (classif) at (15.8,-2.2) {Classe binaire\\boiteux / non};
        \node[small, fill=rulegreen!6] (episode) at (12.2,-4.0) {Épisode\\probable};
        \node[small, fill=rulegreen!6] (alerte) at (15.8,-4.0) {Alerte\\actionnable};

        \draw[arrow] (data) -- (model);
        \draw[arrow] (model) -- (output);

        % Lignes: parent -> rang 1 (directes, courtes)
        \draw[black!30] (data.south) -- ++(0,-0.4) -| (obs.north);
        \draw[black!30] (data.south) -- ++(0,-0.4) -| (vision.north);
        \draw[black!30] (model.south) -- ++(0,-0.4) -| (rules.north);
        \draw[black!30] (model.south) -- ++(0,-0.4) -| (sup.north);
        \draw[black!30] (output.south) -- ++(0,-0.4) -| (score.north);
        \draw[black!30] (output.south) -- ++(0,-0.4) -| (classif.north);

        % Lignes: rang 1 -> rang 2 (verticales simples)
        \draw[black!30] (obs.south) -- (accel.north);
        \draw[black!30] (vision.south) -- (multi.north);
        \draw[black!30] (rules.south) -- (unsup.north);
        \draw[black!30] (sup.south) -- (hyb.north);
        \draw[black!30] (score.south) -- (episode.north);
        \draw[black!30] (classif.south) -- (alerte.north);

        % Encadrement du positionnement de ce mémoire (autour des nœuds unsup et hyb)
        \node[draw=rulegreen, thick, rounded corners=3pt, inner sep=6pt,
              fit=(unsup)(hyb)] (encadre) {};
        \node[font=\footnotesize\bfseries, rulegreen!80!black, below=4pt of encadre]
            {Positionnement de ce mémoire};
    \end{tikzpicture}%
    }
    \caption{Taxonomie simplifiée des approches de détection automatisée de boiterie.}
    \label{fig:taxonomie_approches}
\end{figure}

Du point de vue de la stratégie analytique, les méthodes de détection se déclinent
classiquement en deux familles:
\begin{itemize}
    \item[(a)] \textit{Modèles supervisés}: lorsque des labels fiables et suffisamment
    représentatifs sont disponibles. Ces approches peuvent atteindre de bonnes performances
    de classification
    \citep{vanhertem2014evaluation,viazzi2013analysis,blackie2011automatic}, mais leur
    dépendance aux labels limite leur déployabilité dans des contextes où l'étiquetage est
    incomplet ou incertain;
    \item[(b)] \textit{Détection d'anomalies et règles expertes}: lorsque les labels sont
    partiels ou absents. Ces approches ne nécessitent pas d'exemples positifs pour
    l'entraînement et sont donc mieux adaptées aux situations réelles où la boiterie n'est
    pas systématiquement diagnostiquée et enregistrée.
\end{itemize}

Dans un contexte applicatif réel, la supervision exhaustive est rarement disponible.
Cette contrainte favorise les architectures hybrides combinant score algorithmique
et règles métier interprétables \citep{chandola2009anomaly,palma2025scoping}. Une telle
orientation rejoint des travaux récents de comparaison d'algorithmes non supervisés sur
données bovines, qui montrent que le choix de l'algorithme dépend du compromis visé entre
sensibilité, stabilité et volume de faux positifs, plutôt que d'une supériorité universelle
d'une méthode unique \citep{michelena2025comparative}.

\subsection{Isolation Forest}
L'algorithme Isolation Forest (IF), proposé par Liu et al. \citep{liu2008isolation}, repose
sur le principe que les anomalies sont plus faciles à isoler que les points normaux dans un
espace de \textit{features}. Sa complexité quasi-linéaire en
$O(n \log n)$\footnote{La notation $O(\cdot)$ (dite « grand~O ») décrit la façon dont le
coût de calcul croît avec le nombre d'observations~$n$~: $O(n \log n)$ correspond à un
coût quasi proportionnel à~$n$, alors que $O(n^2)$ croît comme le carré de~$n$ et devient
rapidement prohibitif pour de grands volumes de données.} le rend particulièrement
adapté aux applications où le volume de données croît avec la taille du troupeau. IF est
largement utilisé dans les applications de détection d'anomalies en série temporelle, et
son implémentation dans scikit-learn (bibliothèque Python de référence pour l'apprentissage
automatique) \citep{sklearn_api} en fait un choix pratique pour les systèmes de production.

La Figure~\ref{fig:if_principle} illustre le principe fondamental d'IF: une anomalie
(point rouge) est isolée en seulement deux coupures aléatoires, tandis que les points
normaux (cluster bleu) nécessitent de nombreuses partitions supplémentaires. Le score
d'anomalie est inversement proportionnel au nombre de coupures.

\begin{figure}[H]
    \centering
    \resizebox{0.65\textwidth}{!}{%
    \begin{tikzpicture}
        % Cadre
        \draw[black!50] (0,0) rectangle (12,6);

        % --- Cluster normal : 15 points groupés à gauche ---
        \foreach \x/\y in {2.0/3.0, 2.3/3.4, 2.6/2.8, 2.2/3.2, 1.9/3.5,
                           2.5/3.0, 2.1/2.7, 2.4/3.3, 2.7/3.5, 1.8/3.1,
                           2.2/2.9, 2.5/3.4, 2.3/2.8, 2.1/3.3, 2.6/3.2} {
            \fill[ifblue] (\x,\y) circle (2.5pt);
        }

        % --- Anomalie : 1 point isolé à droite ---
        \fill[rawred] (9.5,4.5) circle (4pt);

        % --- Coupure 1 : verticale entre cluster et anomalie ---
        \draw[black!70, dashed, thick] (6,0.3) -- (6,5.7);

        % --- Coupure 2 : horizontale, zone droite seulement ---
        \draw[black!70, dashed, thick] (6.3,2.8) -- (11.7,2.8);

        % --- Étiquettes hors figure (marges) ---
        % Cluster
        \node[font=\footnotesize, black] at (2.3,1.5) {Points normaux};
        \node[font=\scriptsize, black!50] at (2.3,1.0) {difficiles à isoler};

        % Anomalie
        \node[font=\footnotesize, black] at (9.5,5.3) {Anomalie};
        \node[font=\scriptsize, black!50] at (9.5,3.7) {isolée en 2 coupures};

        % Numéros des coupures
        \node[font=\scriptsize, black, fill=white, inner sep=2pt] at (6,0.9) {1};
        \node[font=\scriptsize, black, fill=white, inner sep=2pt] at (10.8,2.8) {2};
    \end{tikzpicture}%
    }
    \caption{Principe d'isolation des anomalies dans Isolation Forest.}
    \label{fig:if_principle}
\end{figure}

\subsection{Local Outlier Factor}
Le Local Outlier Factor (LOF), proposé par Breunig et al. \citep{breunig2000lof}, adopte
une approche différente en estimant la densité locale de chaque point par rapport à ses
voisins. Contrairement à IF qui mesure l'isolabilité globale, LOF détecte les anomalies
locales, c'est-à-dire les points qui sont significativement moins denses que leur voisinage.
Sa complexité en $O(n^2)$ dans le cas général le rend moins apte au passage à l'échelle,
mais il peut capturer des structures que les méthodes globales manquent. Des travaux récents sur données bovines
confirment que LOF présente une sensibilité accrue aux anomalies de densité locale, ce qui
peut être avantageux pour les perturbations comportementales subtiles
\citep{michelena2025comparative}.

\subsection{Comparaison synthétique IF et LOF}
Le Tableau~\ref{tab:if_vs_lof} synthétise les différences fondamentales entre les deux
détecteurs utilisés dans ce projet.

\begin{table}[H]
    \centering
    \caption{Comparaison des propriétés d'Isolation Forest et de Local Outlier Factor.}
    \label{tab:if_vs_lof}
    \begin{tabular}{lll}
        \toprule
        \textbf{Propriété}            & \textbf{Isolation Forest}
            & \textbf{Local Outlier Factor} \\
        \midrule
        Principe                      & Isolabilité par partitions
            & Densité locale relative \\
        Type d'anomalie ciblée        & Globale (isolée dans l'espace)
            & Locale (faible densité relative) \\
        Complexité entraînement       & $O(n \log n)$
            & $O(n^2)$ (cas général) \\
        Complexité prédiction         & $O(n)$
            & $O(n \cdot k)$ \\
        Sensibilité aux cas faibles   & Modérée
            & Élevée \\
        Robustesse bruit ambiant      & Élevée
            & Modérée \\
        Paramètres critiques          & \texttt{contamination}, \texttt{n\_estimators}
            & \texttt{n\_neighbors}, \texttt{contamination} \\
        Scalabilité troupeau          & Compatible ($>$200 vaches)
            & Limitée ($<$50 vaches) \\
        \bottomrule
    \end{tabular}
\end{table}

La comparaison éclaire le choix d'IF comme détecteur principal dans ce projet, car sa
complexité quasi-linéaire et sa robustesse au bruit ambiant en font un candidat naturel
pour un déploiement sur des troupeaux de taille variable. LOF est néanmoins conservé comme point de
comparaison dans l'étude d'ablation (Chapitre~\ref{chapitre:validation}), où sa sensibilité
accrue aux cas faibles pourrait se révéler avantageuse sur certains profils d'injection.

\subsection{Comparaison systématique des méthodes de détection d'anomalies}

Au-delà de la comparaison binaire entre Isolation Forest et LOF, il est utile de situer ces
deux algorithmes dans le paysage plus large des méthodes de détection d'anomalies et de
classification applicables au domaine. Le Tableau~\ref{tab:comparaison_methodes_anomalies}
propose une comparaison systématique de dix approches, couvrant des méthodes non supervisées,
supervisées et semi-supervisées. Les critères y sont évalués pour un déploiement en
élevage de précision; les notations sont les suivantes~: $n$~nombre d'observations,
$d$~nombre de \textit{features}, $h$~nombre de neurones cachés, $k$~nombre de voisins.
Cette mise en perspective permet de justifier de manière plus rigoureuse le choix de
l'algorithme principal retenu dans ce mémoire.

Les critères retenus pour la comparaison sont les suivants:
\begin{itemize}
    \item la complexité algorithmique (notation~$O$);
    \item le caractère supervisé ou non;
    \item le niveau d'interprétabilité de la sortie;
    \item la capacité native de gestion de données multimodales, c'est-à-dire des
    \textit{features} de natures hétérogènes;
    \item une référence fondatrice pour chaque méthode.
\end{itemize}

Une revue comparative récente de \cite{goldstein2016comparative} sur données multivariées
confirme que le choix d'un détecteur d'anomalies dépend fortement du compromis entre
scalabilité, robustesse au bruit et type de structure recherchée.

\begin{table}[H]
    \centering
    \small
    \setlength{\tabcolsep}{2.5pt}
    \caption{Comparaison systématique de dix méthodes de détection d'anomalies et de
    classification.}
    \label{tab:comparaison_methodes_anomalies}
     \begin{tabular}{lccccl}
        \toprule
        \textbf{Méthode} & \textbf{Compl.} & \textbf{Supervisé} &
        \textbf{Interprétable} & \textbf{Multimodal} &
        \textbf{Réf. clé} \\
        \midrule
        Isolation Forest
        & $O(n \log n)$ & Non & Partiel & Oui
        & \cite{liu2008isolation} \\
        \addlinespace
        Local Outlier Factor
        & $O(n^2)$ & Non & Partiel & Oui
        & \cite{breunig2000lof} \\
        \addlinespace
        One-Class SVM
        & $O(n^2$--$n^3)$ & Non & Non & Oui
        & \cite{scholkopf2001estimating} \\
        \addlinespace
        Autoencoder
        & $O(n \cdot d \cdot h)$ & Non & Non & Oui
        & \cite{sakurada2014anomaly} \\
        \addlinespace
        DBSCAN
        & $O(n \log n)$ & Non & Partiel & Oui
        & \cite{ester1996density} \\
        \addlinespace
        $k$-NN anomaly
        & $O(n^2)$ & Non & Oui & Non
        & \cite{ramaswamy2000efficient} \\
        \addlinespace
        Random Forest (supervisé)
        & $O(n \log n)$ & Oui & Oui & Oui
        & \cite{breiman2001} \\
        \addlinespace
        Régression logistique
        & $O(n \cdot d)$ & Oui & Oui & Oui
        & \cite{cox1958} \\
        \addlinespace
        LSTM / RNN
        & $O(n \cdot h^2)$ & Oui & Non & Oui
        & \cite{hochreiter1997} \\
        \addlinespace
        Elliptic Envelope
        & $O(n \cdot d^2)$ & Non & Oui & Oui
        & \cite{rousseeuw1999fast} \\
        \bottomrule
    \end{tabular}
\end{table}

Les autres méthodes du tableau se répartissent en familles complémentaires. Les approches
par densité (DBSCAN, LOF) repèrent les points situés dans des zones peu denses de l'espace~;
les méthodes à frontière (One-Class SVM, Elliptic Envelope) apprennent une enveloppe du
comportement normal et signalent ce qui en sort~; les approches neuronales (Autoencoder,
LSTM) modélisent la normalité par reconstruction et détectent les écarts de reconstruction~;
les méthodes par distance ($k$-NN) mesurent l'éloignement aux plus proches voisins~; enfin,
les modèles supervisés (Random Forest, régression logistique) exigent des labels
d'entraînement. Isolation Forest se distingue de l'ensemble par son principe d'isolation
directe des anomalies, sans hypothèse préalable de densité ni de distance.

Plusieurs constats se dégagent de cette comparaison. Parmi les méthodes non supervisées, seuls
Isolation Forest et DBSCAN offrent une complexité quasi-linéaire en $O(n \log n)$, ce qui
les rend compatibles avec un passage à l'échelle sur des troupeaux de grande taille. Les
approches supervisées (Random Forest, régression logistique, LSTM) peuvent atteindre de bonnes
performances de classification, mais elles nécessitent des labels fiables et suffisamment
représentatifs, une condition rarement satisfaite dans les contextes d'élevage réel
\citep{alsaaod2019automatic,oleary2020accelerometers}.

Le choix d'Isolation Forest comme détecteur principal dans ce mémoire repose sur la
convergence de cinq arguments:
\begin{enumerate}
    \item \textit{Fonctionnement non supervisé}: IF ne nécessite pas de labels positifs
    (cas de boiterie confirmés) pour l'entraînement, ce qui le rend directement applicable
    dans des contextes où le diagnostic clinique n'est pas systématiquement enregistré;
    \item \textit{Scalabilité}: sa complexité en $O(n \log n)$ permet de traiter des volumes
    de données croissants sans dégradation significative du temps de calcul, contrairement
    aux approches en $O(n^2)$ comme LOF ou $k$-NN;
    \item \textit{Interprétabilité partielle}: le score d'anomalie produit par IF est
    directement exploitable comme indicateur de sévérité et peut être combiné avec des
    seuils métier pour produire des alertes graduées et compréhensibles;
    \item \textit{Performances établies sur données tabulaires}: la littérature démontre
    l'efficacité d'IF sur des jeux de données tabulaires multivariés, un profil cohérent
    avec les \textit{features} agrégées extraites de capteurs comportementaux
    \citep{liu2012isolation_extended,goldstein2016comparative};
    \item \textit{Support natif dans scikit-learn}: l'implémentation de référence dans
    scikit-learn \citep{sklearn_api} garantit une intégration directe dans un pipeline
    Python de production, sans dépendance à des frameworks spécialisés.
\end{enumerate}

\subsection{Protocole de sélection des études}

La présente revue couvre les travaux portant sur la détection automatisée de boiterie
bovine par capteurs, publiés entre 1990 et 2025. La sélection a été conduite par
consultation des bases Google~Scholar, Web~of~Science et PubMed, avec les termes clés
\textit{dairy cattle lameness detection}, \textit{bovine lameness accelerometer},
\textit{anomaly detection livestock} et leurs équivalents francophones.

\textit{Critères d'inclusion :}
\begin{itemize}
    \item Études portant sur la détection ou la prédiction de la boiterie chez les bovins
    laitiers, y compris les revues de littérature qui synthétisent des données empiriques;
    \item Utilisant des données de capteurs comportementaux (accéléromètres, podomètres,
    tapis de pesée, systèmes de traite automatisés) ou de vision par ordinateur;
    \item Rapportant au moins un indicateur de performance quantitatif (sensibilité,
    spécificité, AUC, taux de faux positifs) ou présentant une comparaison systématique
    de méthodes;
    \item Publiées en anglais ou en français, dans des revues à comité de lecture ou des
    actes de conférences indexées.
\end{itemize}

\textit{Critères d'exclusion :}
\begin{itemize}
    \item Études portant exclusivement sur des espèces autres que les bovins laitiers;
    \item Articles sans validation empirique (éditoriaux, perspectives théoriques sans
    données);
    \item Travaux consacrés uniquement à l'étiologie, au traitement ou à la chirurgie
    podologique, sans composante de détection automatisée.
\end{itemize}

Les études retenues pour la synthèse quantitative
(Tableau~\ref{tab:performances_litterature}) ont été sélectionnées pour leur
représentativité des grandes familles d'approches identifiées dans la taxonomie
(Figure~\ref{fig:taxonomie_approches}). Le présent mémoire se situe dans la branche
capteurs comportementaux~$\rightarrow$ détection d'anomalies~$\rightarrow$ règles
métier~$\rightarrow$ alerte actionnable. La revue ne prétend pas à l'exhaustivité d'une
méta-analyse systématique: elle vise à couvrir les travaux qui ont le plus directement
influencé les décisions méthodologiques du présent projet, et à mettre en évidence la
variabilité des protocoles qui complique les comparaisons directes de performance.

% =====================================================================
\subsection{Travaux représentatifs et implications pour ce mémoire}
Le Tableau~\ref{tab:revue_comparee} résume quelques travaux représentatifs couvrant des
angles différents de la détection de boiterie: revues méthodologiques, études multi-fermes,
apprentissage profond sur signaux accélérométriques et comparaisons d'algorithmes non
supervisés. L'objectif n'est pas de dresser un inventaire exhaustif, mais de mettre en
évidence les lignes de force qui ont orienté les choix méthodologiques de ce mémoire.

\begin{table}[H]
    \centering
    \small
    \caption{Synthèse de travaux représentatifs sur la détection automatisée de boiterie
    bovine.}
    \label{tab:revue_comparee}
    \begin{tabular}{>{\raggedright\arraybackslash}p{2.8cm}
                    >{\raggedright\arraybackslash}p{2.4cm}
                    >{\raggedright\arraybackslash}p{2.7cm}
                    >{\raggedright\arraybackslash}p{3.1cm}
                    >{\raggedright\arraybackslash}p{3.2cm}}
        \toprule
        \textbf{Travail} & \textbf{Données} & \textbf{Approche} & \textbf{Apport principal} & \textbf{Leçon pour ce mémoire} \\
        \midrule
        \citet{oleary2020accelerometers}
        & Revues d'accélérométrie
        & Revue ciblée
        & Accéléromètres utiles, mais protocoles très variables
        & Documenter capteurs, labels et métriques \\
        \addlinespace
        \citet{alsaaod2019automatic}
        & Revue de systèmes automatiques
        & Revue de littérature
        & Pas d'alerte précoce validée de bout en bout
        & Viser un système complet, pas un modèle isolé \\
        \addlinespace
        \citet{lavrova2023lameness}
        & Six fermes ; activité, couchage, production
        & Modèle logistique à effets mixtes
        & Sensibilité accrue en combinant comportement et contexte
        & Privilégier des \textit{features} agrégées hétérogènes \\
        \addlinespace
        \citet{balasso2023behaviour}
        & Accéléromètres tri-axiaux
        & Apprentissage profond
        & Capte des motifs comportementaux complexes
        & Plus de signal demande plus de labels et de calibration \\
        \addlinespace
        \citet{michelena2025comparative}
        & Données capteurs bovines
        & Comparaison non supervisée
        & Aucun algorithme universellement dominant
        & Évaluer le compromis sensibilité/stabilité/bruit \\
        \addlinespace
        \citet{palma2025scoping}
        & Revue IA, élevage de précision
        & Scoping review
        & Difficultés de déploiement, transfert et adoption
        & Sortie interprétable et traçable pour le terrain \\
        \bottomrule
    \end{tabular}
\end{table}

Deux enseignements structurent particulièrement ce mémoire. D'une part, la qualité du système
ne peut pas être réduite à une seule métrique de détection: un outil utile en ferme doit
concilier sensibilité, robustesse au bruit et lisibilité des alertes. D'autre part, la
valeur scientifique ne réside pas uniquement dans le choix d'un algorithme, mais aussi dans
la capacité à documenter un protocole reproductible, ses hypothèses et ses limites.

Avant de fixer définitivement le positionnement du mémoire, une dernière famille d'approches
mérite d'être située, à savoir l'apprentissage profond, dont l'essor récent en élevage de
précision soulève la question de sa pertinence face aux méthodes classiques pour la
détection de boiterie.

\subsection{Apprentissage profond versus approches classiques}

L'émergence des réseaux de neurones profonds a renouvelé l'intérêt pour la détection
automatisée de boiterie, notamment via les architectures convolutives (CNN) et récurrentes
(LSTM) appliquées à des signaux accélérométriques bruts
\citep{balasso2023behaviour,qiao2021review}. Le
Tableau~\ref{tab:deep_vs_classique} compare les deux familles d'approches sur les dimensions
pertinentes pour un déploiement en élevage, en évaluant les critères dans le contexte
d'un déploiement opérationnel avec labels partiels.

\begin{table}[H]
    \centering
    \caption{Comparaison des approches d'apprentissage profond et classiques pour la détection
    de boiterie bovine.}
    \label{tab:deep_vs_classique}
    \begin{tabular}{>{\raggedright\arraybackslash}p{3.2cm}
                    >{\raggedright\arraybackslash}p{4.5cm}
                    >{\raggedright\arraybackslash}p{4.5cm}}
        \toprule
        \textbf{Critère} & \textbf{Deep Learning (CNN/LSTM)}
        & \textbf{Classique (IF, LOF, seuils)} \\
        \midrule
        Besoin en labels
        & Élevé (milliers d'exemples étiquetés)
        & Faible à nul (aucun label positif requis) \\
        \addlinespace
        Interprétabilité
        & Limitée (\textit{features} apprises peu lisibles)
        & Élevée (z-scores, seuils, règles explicites) \\
        \addlinespace
        Volume de données
        & Important (milliers de séquences par classe)
        & Modeste (centaines de bins par vache) \\
        \addlinespace
        Complexité de déploiement
        & GPU et frameworks lourds (TensorFlow, PyTorch)
        & CPU et bibliothèques légères (scikit-learn) \\
        \addlinespace
        Performances publiées
        & Se 80--92\% avec labels \citep{balasso2023behaviour}
        & Se 76--91\% selon les études \citep{alsaaod2012electronic,chapinal2010measurement} \\
        \addlinespace
        Adaptabilité inter-fermes
        & Ré-entraînement ou transfert requis
        & Adaptation native (modèle par vache) \\
        \bottomrule
    \end{tabular}
\end{table}

Dans le contexte de ce mémoire, trois arguments justifient le choix d'une approche classique
plutôt que profonde:
\begin{enumerate}
    \item \textit{Absence de supervision}: le corpus ne contient pas d'évaluation
    locomotrice synchronisée, ce qui exclut les approches supervisées standard. Les travaux de
    \cite{beer2016objective} et \cite{thorup2015lameness} montrent que même les approches
    classiques supervisées nécessitent des campagnes de notation dédiées;
    \item \textit{Exigence d'interprétabilité}: le système doit produire des alertes
    compréhensibles par les éleveurs, avec des raisons explicites (z-scores, persistance,
    cohérence multi-signaux). Les réseaux profonds, même performants, ne fournissent pas
    nativement ce niveau d'explicabilité \citep{rutten2013invited};
    \item \textit{Contrainte de corpus}: avec 28~vaches sur 6~semaines, le volume de données
    est insuffisant pour entraîner un réseau profond sans surapprentissage. Les détecteurs
    d'anomalies comme IF fonctionnent avec des échantillons bien plus modestes
    \citep{liu2008isolation}.
\end{enumerate}

Privilégier les méthodes classiques ne signifie pas que l'apprentissage profond soit
inadapté à la détection de boiterie en général. Les travaux de \cite{balasso2023behaviour}
et de \cite{qiao2021review} démontrent le potentiel de ces méthodes lorsque les conditions
sont réunies (labels, volume, GPU). Cependant, dans le
contexte opérationnel de ce mémoire, l'approche classique hybride est plus pragmatique et
défendable.

% =====================================================================
\section{Technologies capteurs pour la détection de boiterie}
\label{sec:technologies_capteurs}
% =====================================================================

Avant d'aborder l'hétérogénéité des protocoles, cette section poursuit un double objectif:
situer d'abord brièvement les grandes familles de capteurs disponibles pour la détection
de boiterie, puis décrire le capteur effectivement retenu dans ce projet (IceQube) ainsi
que les variables comportementales qu'il produit. Les systèmes de surveillance
comportementale reposent sur plusieurs familles de capteurs, dont les caractéristiques
conditionnent directement les variables disponibles et les performances atteignables.

Le Tableau~\ref{tab:technologies_capteurs} présente les principales technologies et leurs
caractéristiques opérationnelles. On distingue trois grandes familles: les capteurs portés
par l'animal, tels que les accéléromètres, les podomètres et les colliers; les capteurs
intégrés à l'infrastructure, comme les tapis de pesée et les systèmes de traite; et les
systèmes sans contact, comme les caméras et la vision par ordinateur. Les accéléromètres
portés constituent la technologie la plus étudiée dans la littérature récente
\citep{oleary2020accelerometers,alsaaod2019automatic}.

\begin{table}[H]
    \centering
    \small
    \caption{Technologies capteurs utilisées dans la détection de boiterie bovine.}
    \label{tab:technologies_capteurs}
    \begin{tabular}{>{\raggedright\arraybackslash}p{2.4cm}
                    >{\raggedright\arraybackslash}p{2.8cm}
                    >{\raggedright\arraybackslash}p{2.8cm}
                    >{\raggedright\arraybackslash}p{2.2cm}
                    >{\raggedright\arraybackslash}p{2.2cm}}
        \toprule
        \textbf{Technologie} & \textbf{Variables mesurées} & \textbf{Avantages} &
        \textbf{Limites} & \textbf{Études repr.} \\
        \midrule
        Accéléromètre (patte/collier)
        & Activité, pas, transitions couché/debout, durée couchage
        & Continu, individuel, coût modéré
        & Dérive, position variable
        & \cite{alsaaod2012electronic,balasso2023behaviour} \\
        \addlinespace
        Podomètre
        & Nombre de pas, distance parcourue
        & Simple, robuste
        & Signal limité, faible résolution
        & \cite{flower2006effect,mazrier2006reliability} \\
        \addlinespace
        Tapis de pesée
        & Distribution du poids entre les pieds
        & Objectif, quantitatif
        & Infrastructure fixe, coût élevé
        & \cite{alsaaod2012electronic,hempel2019dairy} \\
        \addlinespace
        Vision par ordinateur
        & Posture, arc du dos, démarche
        & Sans contact, riche en information
        & Sensible aux conditions, calibration
        & \cite{qiao2021review} \\
        \addlinespace
        Système de traite automatisé
        & Production laitière, conductivité, fréquentation
        & Déjà installé, données continues
        & Signal indirect pour la boiterie
        & \cite{lavrova2023lameness,vanhertem2013lameness} \\
        \bottomrule
    \end{tabular}
\end{table}

Dans le cadre de ce mémoire, les données proviennent de capteurs IceQube
(IceRobotics Ltd, Édimbourg, Royaume-Uni), des accéléromètres tri-axiaux conçus
spécifiquement pour le suivi comportemental des vaches laitières. La
Figure~\ref{fig:iceqube_capteur} montre la représentation schématique du capteur IceQube,
dont le boîtier étanche contient l'accéléromètre tri-axial et se fixe à la patte par un
bracelet en plastique muni d'une boucle et de trous d'ajustement. La
Figure~\ref{fig:iceqube_position} illustre son positionnement sur le métatarse de la patte
arrière~; les axes $x$, $y$, $z$ de l'accéléromètre permettent de distinguer la position
debout (patte verticale, $\vec{g}$ aligné avec $z$) de la position couchée (patte
horizontale).

\begin{figure}[H]
    \centering
    \resizebox{0.7\textwidth}{!}{%
    \begin{tikzpicture}[>=Stealth]
        % --- Bracelet gauche (vers la boucle) ---
        \fill[ifblue!35, rounded corners=2pt]
            (-3.0,0.65) rectangle (0,1.35);
        % Boucle de fermeture
        \fill[ifblue!50, rounded corners=2pt]
            (-3.6,0.5) rectangle (-3.0,1.5);
        \draw[ifblue!70, thick, rounded corners=1pt]
            (-3.45,0.7) rectangle (-3.15,1.3);

        % --- Bracelet droit (trous) ---
        \fill[ifblue!35, rounded corners=2pt]
            (4.0,0.65) rectangle (7.0,1.35);
        % Trous d'ajustement (ovales)
        \foreach \x in {4.6, 5.2, 5.8, 6.4} {
            \fill[white] (\x,1.0) ellipse (0.12 and 0.15);
        }

        % --- Boîtier central (vue de face) ---
        \fill[ifblue!50, rounded corners=4pt]
            (0,0) rectangle (4.0,2.0);
        % Bord supérieur légèrement plus clair
        \fill[ifblue!38, rounded corners=3pt]
            (0.1,1.55) rectangle (3.9,1.95);

        % Étiquette blanche
        \fill[white, rounded corners=2pt]
            (0.6,0.35) rectangle (3.4,1.4);
        \node[font=\small\bfseries, black] at (2.0,1.05) {IceQube};
        \node[font=\scriptsize, black!60] at (2.0,0.7) {IceRobotics};

        % --- Annotations ---
        \draw[->, black!55, thick]
            (4.5,2.5) -- (3.5,2.05)
            node[pos=0, anchor=south west, font=\scriptsize, black!70]
            {Boîtier étanche};

        \draw[->, black!55, thick]
            (-3.6,0.0) -- (-3.45,0.45)
            node[pos=0, anchor=north, font=\scriptsize, black!70]
            {Boucle};

        \draw[->, black!55, thick]
            (7.0,0.0) -- (5.8,0.6)
            node[pos=0, anchor=north, font=\scriptsize, black!70]
            {Trous d'ajustement};
    \end{tikzpicture}%
    }%
    \caption{Représentation schématique du capteur IceQube.}
    \label{fig:iceqube_capteur}
\end{figure}

\begin{figure}[H]
    \centering
    \resizebox{0.4\textwidth}{!}{%
    \begin{tikzpicture}[>=Stealth]
        % --- Patte (forme simple, allongée verticalement) ---
        \fill[brown!12, rounded corners=5pt]
            (0,0) -- (0,6.0) -- (0.4,6.5) -- (1.4,6.5) -- (1.8,6.0)
            -- (1.8,0) -- (1.5,-0.4) -- (0.3,-0.4) -- cycle;
        \draw[brown!30, thick, rounded corners=5pt]
            (0,0) -- (0,6.0) -- (0.4,6.5) -- (1.4,6.5) -- (1.8,6.0)
            -- (1.8,0) -- (1.5,-0.4) -- (0.3,-0.4) -- cycle;

        % Jarret (ligne courte, texte à droite)
        \draw[brown!35, thick] (-0.2,4.5) -- (2.2,4.5);
        \node[font=\scriptsize, brown!55, anchor=west] at (2.4,4.5) {Jarret};

        % Zone métatarse (lignes courtes à droite seulement)
        \draw[brown!25, dashed] (1.9,1.0) -- (2.8,1.0);
        \draw[brown!25, dashed] (1.9,3.2) -- (2.8,3.2);
        \draw[decorate, decoration={brace, amplitude=4pt, mirror}, brown!40]
            (2.9,1.0) -- (2.9,3.2)
            node[midway, right=5pt, font=\scriptsize, brown!55] {Métatarse};

        % --- Bracelet autour de la patte ---
        \fill[ifblue!35, rounded corners=2pt]
            (-0.6,1.7) rectangle (0,2.5);
        \fill[ifblue!35, rounded corners=2pt]
            (1.8,1.7) rectangle (2.4,2.5);
        % Boucle côté gauche
        \fill[ifblue!50, rounded corners=1pt]
            (-0.8,1.8) rectangle (-0.6,2.4);

        % --- Capteur (sur la face avant de la patte) ---
        \fill[ifblue!50, rounded corners=3pt]
            (0.0,1.6) rectangle (1.8,2.6);
        \fill[white, rounded corners=1pt]
            (0.25,1.75) rectangle (1.55,2.45);
        \node[font=\scriptsize\bfseries, black] at (0.9,2.1) {IceQube};

        % --- Axes de l'accéléromètre (en bas à gauche, bien isolés) ---
        \coordinate (AX) at (-2.2,0.0);

        % z vers le haut
        \draw[->, rawred, very thick] (AX) -- ++(0,1.4)
            node[above, font=\small, rawred] {$z$};

        % x vers la droite
        \draw[->, rulegreen, very thick] (AX) -- ++(1.2,0)
            node[right, font=\small, rulegreen] {$x$};

        % y vers l'avant (hors plan)
        \draw[->, ifblue!80!black, very thick] (AX) -- ++(-0.7,-0.7)
            node[below left, font=\small, ifblue!80!black] {$y$};

        % --- Gravité (à droite, au-dessus du jarret) ---
        \draw[->, black!45, thick, dashed] (3.8,6.0) -- ++(0,-1.3)
            node[midway, right=2pt, font=\scriptsize, black!50] {$\vec{g}$};
    \end{tikzpicture}%
    }%
    \caption{Positionnement du capteur IceQube sur la patte arrière.}
    \label{fig:iceqube_position}
\end{figure}

Le capteur se fixe à la patte arrière de l'animal (métatarse) par un bracelet à boucle, sans
gêner la locomotion ni le couchage. L'accéléromètre interne mesure en continu l'accélération
sur trois axes. La composante gravitationnelle permet de distinguer la position couchée
(patte horizontale) de la position debout (patte verticale), tandis que les variations
dynamiques d'accélération quantifient l'activité locomotrice. Un algorithme embarqué agrège
ces mesures brutes en intervalles configurables (dans ce projet, des bins de 15~minutes),
un choix validé par \cite{ledgerwood2010lying} comme offrant un bon compromis entre
résolution temporelle et stabilité des mesures de couchage. La fiabilité de cette approche
est documentée par plusieurs études indépendantes: \cite{alsaaod2012electronic} rapportent
une forte concordance entre les mesures d'accéléromètres de patte et l'observation directe
pour la classification couché/debout, tandis que \cite{blackie2011automatic} confirment la
fiabilité des accéléromètres pour le suivi longitudinal de l'activité.

Chaque bin de 15~minutes génère sept colonnes de données, dont cinq sont utilisées dans le
pipeline:
\begin{itemize}
    \item \textit{Steps}: nombre de pas détectés par l'accéléromètre;
    \item \textit{Motion Index}: indice synthétique d'activité locomotrice (unité
    propriétaire IceRobotics, dérivée des variations d'accélération);
    \item \textit{Lying Time}: durée cumulée en position couchée (en minutes, bornée à 15);
    \item \textit{Standing Time}: durée cumulée en position debout (en minutes, bornée à 15);
    \item \textit{Transitions}: nombre total de changements de posture
    couché$\leftrightarrow$debout.
\end{itemize}
Les colonnes supplémentaires (\textit{Transitions Up} et \textit{Transitions Down})
détaillent la direction des transitions mais ne sont pas exploitées directement par le
pipeline, qui utilise leur somme.

Les données agrégées sont stockées dans la mémoire interne du capteur, puis récupérées
périodiquement. Elles sont accessibles via la plateforme web CowAlert (IceRobotics), qui
offre une interface de visualisation en ligne, et peuvent être exportées au format tabulaire
(CSV), directement exploitable par le pipeline de traitement décrit dans ce mémoire.

Le modèle IceTag, prédécesseur de l'IceQube, fournit des variables similaires (pas, temps
couché, transitions). Les deux modèles sont largement utilisés dans la recherche sur la
boiterie bovine \citep{beer2016objective,alsaaod2012electronic}. Le choix de l'IceQube
dans ce projet s'inscrit dans la tendance dominante de la littérature, où les
accéléromètres de patte représentent la technologie la plus étudiée pour la détection
automatisée de boiterie \citep{oleary2020accelerometers,bewley2008review}.

\cite{ito2010lying} montrent que les vaches boiteuses modifient significativement leur
comportement de couchage (durée allongée, fréquence de transitions réduite), ce qui
justifie l'inclusion de ces variables dans le pipeline. \cite{nechanitzky2016analysis}
confirment que les changements comportementaux précèdent souvent les signes cliniques
visibles, renforçant l'intérêt d'une détection précoce par capteurs.

% =====================================================================
\section{Hétérogénéité des protocoles~: un obstacle à la comparaison}
% =====================================================================

Les revues spécialisées sur la détection automatisée de boiterie convergent vers une
limite récurrente: l'hétérogénéité des protocoles. Le type de capteur, sa position sur
l'animal, la définition opérationnelle de la boiterie, la taille des cohortes et l'horizon
d'évaluation varient fortement d'une étude à l'autre, ce qui complique la comparaison
directe des performances publiées, comme le soulignent les travaux de
\citet{oleary2020accelerometers} et de \citet{alsaaod2019automatic}.

Le Tableau~\ref{tab:performances_litterature} présente une synthèse quantitative de seize
travaux représentatifs couvrant la période 1997--2026 (Se~=~sensibilité, Sp~=~spécificité;
«~\textnormal{—}~»~indique une étude descriptive sans métrique de classification,
«~\textnormal{n.d.}~»~une valeur non disponible dans l'article). Les sensibilités publiées
varient de 22\% à 91\%, reflétant l'hétérogénéité des capteurs, des populations et des
définitions de cas. Cette dispersion illustre l'impossibilité de comparer directement les
performances sans tenir compte du protocole de validation, un point développé dans la
section suivante.

\nocite{sprecher1997lameness, flower2006effect, chapinal2010measurement,
alsaaod2012electronic, vanhertem2013lameness, thorup2015lameness, beer2016objective,
nechanitzky2016analysis, hempel2019dairy, oleary2020accelerometers, qiao2021review,
gonzalezgarcia2023systematic, lavrova2023lameness, balasso2023behaviour,
michelena2025comparative}
{\small
\renewcommand{\arraystretch}{1.0}%
\setlength{\tabcolsep}{3pt}%
\tolerance=9999%
\emergencystretch=5em%
\hbadness=10000%
\hfuzz=10pt%
\begin{longtable}{>{\raggedright\arraybackslash}p{3.2cm}
                >{\raggedright\arraybackslash}p{2.3cm}
                >{\raggedright\arraybackslash}p{3.0cm}
                >{\centering\arraybackslash}p{1.0cm}
                >{\centering\arraybackslash}p{1.3cm}
                >{\centering\arraybackslash}p{1.0cm}
                >{\raggedright\arraybackslash}p{2.2cm}}
\caption[Synthèse quantitative de travaux représentatifs en détection de boiterie
bovine]{Synthèse quantitative de travaux représentatifs sur la détection automatisée
de boiterie bovine.}\label{tab:performances_litterature}\\
\toprule
\textbf{Étude} & \textbf{Capteur} & \textbf{Méthode} &
\textbf{$n$} & \textbf{Se (\%)} & \textbf{Sp (\%)} & \textbf{Validation} \\
\midrule
\endfirsthead

\multicolumn{7}{l}{\textit{Tableau \thetable{} (suite)}}\\
\toprule
\textbf{Étude} & \textbf{Capteur} & \textbf{Méthode} &
\textbf{$n$} & \textbf{Se (\%)} & \textbf{Sp (\%)} & \textbf{Validation} \\
\midrule
\endhead

\bottomrule
\endlastfoot
        Sprecher et al. (1997) & Observation & Notation locomotrice & — & — & — & Conceptuelle \\
        \addlinespace
        Flower \& Weary (2006) & Observation & Analyse comportementale & 40 & — & — & Descriptive \\
        \addlinespace
        Chapinal et al. (2010) & Accél.~patte & Régression logistique & 48 & 76 & 83 & 1 ferme \\
        \addlinespace
        Alsaaod et al. (2012) & Pédomètre & Seuils sur activité & 36 & 76--91 & 75--88 & 1 ferme \\
        \addlinespace
        Van Hertem et al. (2014) & Traite autom. & SVM + activité & — & n.d. & n.d. & Terrain \\
        \addlinespace
        Thorup et al. (2015) & Accél.~patte (IceTag) & Modèle linéaire mixte & 4 fermes & 86 & 89 & Multi-fermes \\
        \addlinespace
        Beer et al. (2016) & Accél.~patte (IceTag) & Régression + seuils & 150 & 81 & 87 & Multi-fermes \\
        \addlinespace
        Nechanitzky et al. (2016) & Accél.~patte & Analyse de variance & ${\sim}50$ & — & — & Descriptive \\
        \addlinespace
        Hempel et al. (2019) & Tapis de pesée & RF / ANN & ${\sim}50$ & n.d. & n.d. & Terrain \\
        \addlinespace
        O'Leary et al. (2020) & Accél. (multi) & Revue de 15 études & $>$400 & 22--91 & 56--95 & Revue \\
        \addlinespace
        Qiao et al. (2021) & Vision / multi & Revue & $>$500 & Var. & Var. & Revue \\
        \addlinespace
        González-García et al. (2023) & Multi-sources & Revue systématique & $>$1000 & Var. & Var. & Revue syst. \\
        \addlinespace
        Lavrova et al. (2023) & Collier + traite & Logistique mixte & 400+ & 77 & 82 & 6 fermes \\
        \addlinespace
        Balasso et al. (2023) & Accél.~tri-axial & Apprentissage profond & 120 & 85 & 79 & Multi-fermes \\
        \addlinespace
        Michelena et al. (2025) & Capteurs dérivés & Comparaison non supervisée & 50+ & Var. & Var. & Comparaison \\
        \addlinespace
        \textbf{Ce mémoire} (2026) & IceQube (patte) & IF + règles métier & 28 & 90--100\textsuperscript{a} & n.d. & Synthétique \\
\end{longtable}}

\vspace{-0.5em}
\noindent\textit{\textsuperscript{a}Sur cas forts injectés~; validation synthétique, non
clinique (voir Chapitre~\ref{chapitre:validation}).
Van~Hertem et al.~\citep{vanhertem2013lameness} et Hempel et al.~\citep{hempel2019dairy}
rapportent principalement des corrélations avec scores de locomotion plutôt que des
métriques Se/Sp directes.}

La revue d'O'Leary et al. \citep{oleary2020accelerometers} sur les accéléromètres pour la
détection de boiterie confirme cette dispersion et souligne l'absence de standardisation
dans les protocoles d'évaluation.
Alsaaod et al. \citep{alsaaod2019automatic} soulignent en particulier l'absence
de véritables études de longue durée ayant abouti à un système d'alerte précoce valide de
bout en bout.

Des travaux plus récents confirment l'intérêt des accéléromètres à l'échelle multi-fermes,
mais montrent aussi que les variables utiles mélangent activité, couchage, production et
facteurs individuels plutôt qu'un unique signal universel. Dans une étude longitudinale sur
six fermes allemandes, Lavrova et al. \citep{lavrova2023lameness} obtiennent une sensibilité
d'environ 77\% avec un modèle logistique à effets mixtes, en combinant données d'activité,
durée des épisodes de couchage, production laitière et contexte animal.

Enfin, la qualité de la mesure amont reste déterminante: l'échantillonnage, l'édition des
épisodes et la représentation des comportements influencent directement la fiabilité des
variables dérivées \citep{ledgerwood2010lying,balasso2023behaviour}.

Concrètement, l'hétérogénéité des protocoles se manifeste sur au moins cinq dimensions:
\begin{enumerate}
    \item \textit{Définition du cas}: classe clinique, score locomoteur seuil ou probabilité
    de boiterie selon les études, d'où des phénomènes mesurés et des sensibilités
    difficilement comparables;
    \item \textit{Instrumentation}: position du capteur (patte
    \citep{alsaaod2012electronic,thorup2015lameness}, collier \citep{blackie2011automatic},
    corps entier \citep{viazzi2013analysis}), fréquence d'échantillonnage et signaux
    contextuels varient, produisant des variables de nature différente;
    \item \textit{Granularité analytique}: prédiction instantanée, journalière, par épisode
    ou par vache, avec des exigences de précision temporelle très différentes;
    \item \textit{Objectif métier}: classification, détection précoce, filtrage ou score de
    risque, chaque objectif déplaçant le compromis entre faux positifs et faux négatifs;
    \item \textit{Évaluation}: sensibilité, spécificité, AUC, taux d'alertes, IoU ou
    robustesse inter-fermes, un même système paraissant excellent ou médiocre selon la
    métrique.
\end{enumerate}

Autrement dit, deux articles peuvent rapporter une ``bonne performance'' tout en répondant à
des problèmes différents. Cette observation justifie le choix de consacrer un chapitre entier
aux protocoles de validation et à leurs artefacts, plutôt que de se limiter à un score unique
agrégé.

La Figure~\ref{fig:evidence_map} synthétise deux lacunes structurelles de la littérature~:
la quasi-totalité des méthodes non supervisées n'a été évaluée que sur une seule ferme, et
aucune ne dispose d'une validation synthétique structurée, ce que le présent mémoire vise à
combler (cellule verte, en bas à droite).

%% ---------------------------------------------------------------
%% Matrice de couverture des travaux représentatifs
%% Lignes = famille de méthode | Colonnes = contexte de validation
%% Couleur cellule = niveau de Se% | Cellule rouge = lacune
%% ---------------------------------------------------------------
\begin{figure}[H]
\centering
\resizebox{\textwidth}{!}{%
\begin{tikzpicture}[
  x=1.92cm, y=1.18cm,
  font=\scriptsize,
  every node/.style={inner sep=1.5pt},
]

  %% ===== EN-TÊTE DES COLONNES =====
  %% Colonnes montées à y=4.95 ; flèche à y=4.55 ; ligne séparatrice à y=4.32
  %% → évite le chevauchement entre la 2e ligne des labels et la flèche
  \node[font=\tiny\bfseries, align=center] at (0.5, 4.95) {Conceptuel};
  \node[font=\tiny\bfseries, align=center] at (1.5, 4.95) {1 ferme};
  \node[font=\tiny\bfseries, align=center] at (2.5, 4.95) {Terrain\\[-2pt]pratique};
  \node[font=\tiny\bfseries, align=center] at (3.5, 4.95) {Multi-\\[-2pt]fermes};
  \node[font=\tiny\bfseries, align=center] at (4.5, 4.95) {Revue\,/\\[-2pt]Méta};
  \node[font=\tiny\bfseries, align=center,
        rulegreen!55!black] at (5.5, 4.95) {Synthé-\\[-2pt]tique};

  \node[font=\tiny, black!65, align=center] at (3.0, 4.55)
      {$\longleftarrow$ Validité externe croissante $\longrightarrow$};
  \draw[gray!28, thin] (0, 4.32) -- (6, 4.32);

  %% ===== EN-TÊTE DES RANGÉES =====
  \node[ifblue!80!black, font=\tiny\bfseries, anchor=east, align=right]
      at (-0.06, 3.5) {Seuils /\\règles};
  \node[loforange!80!black, font=\tiny\bfseries, anchor=east, align=right]
      at (-0.06, 2.5) {Supervisé\\classique};
  \node[rawred!80!black, font=\tiny\bfseries, anchor=east, align=right]
      at (-0.06, 1.5) {App.\\profond};
  \node[rulegreen!80!black, font=\tiny\bfseries, anchor=east, align=right]
      at (-0.06, 0.5) {Non\\supervisé};

  %% ==========================================================
  %% ROW 3 — SEUILS / RÈGLES
  %% ==========================================================

  %% Conceptuel : Sprecher '97, n.d.
  \fill[gray!13] (0,3) rectangle (1,4);
  \draw[black!15] (0,3) rectangle (1,4);
  \node[font=\tiny\bfseries] at (0.5, 3.70) {1 étude};
  \node[font=\tiny, black!65] at (0.5, 3.44) {Sprecher~'97};
  \node[font=\tiny, black!58] at (0.5, 3.20) {\itshape conceptuel};

  %% 1 ferme : Flower + Alsaaod + Nechanitzky — Se ≈ 84 %*
  \fill[rulegreen!20] (1,3) rectangle (2,4);
  \draw[black!18] (1,3) rectangle (2,4);
  \node[font=\tiny\bfseries] at (1.5, 3.72) {3 études};
  \node[font=\tiny] at (1.5, 3.50) {Se\,$\approx$\,84\,\%$^{*}$};
  \node[font=\tiny, black!65] at (1.5, 3.26) {Flower~\emph{et~al.}};

  %% Terrain : vide
  \fill[white] (2,3) rectangle (3,4);
  \draw[gray!20, dashed, thin] (2,3) rectangle (3,4);

  %% Multi-fermes : vide
  \fill[white] (3,3) rectangle (4,4);
  \draw[gray!20, dashed, thin] (3,3) rectangle (4,4);

  %% Revue : O'Leary '20 — Se 22–91 % (forte dispersion)
  \fill[loforange!18] (4,3) rectangle (5,4);
  \draw[black!18] (4,3) rectangle (5,4);
  \node[font=\tiny\bfseries] at (4.5, 3.72) {1 étude};
  \node[font=\tiny] at (4.5, 3.50) {Se\,22–91\,\%};
  \node[font=\tiny, black!60] at (4.5, 3.26) {\itshape forte dispersion};

  %% Synthétique : vide
  \fill[white] (5,3) rectangle (6,4);
  \draw[gray!20, dashed, thin] (5,3) rectangle (6,4);

  %% ==========================================================
  %% ROW 2 — SUPERVISÉ CLASSIQUE
  %% ==========================================================

  %% Conceptuel : vide
  \fill[white] (0,2) rectangle (1,3);
  \draw[gray!20, dashed, thin] (0,2) rectangle (1,3);

  %% 1 ferme : Chapinal seul — Se = 76 %  (Hempel = Terrain, déplacé)
  \fill[loforange!18] (1,2) rectangle (2,3);
  \draw[black!18] (1,2) rectangle (2,3);
  \node[font=\tiny\bfseries] at (1.5, 2.65) {1 étude};
  \node[font=\tiny] at (1.5, 2.45) {Se\,=\,76\,\%};
  \node[font=\tiny, black!65] at (1.5, 2.24) {Chapinal~'10};

  %% Terrain : van Hertem '14 + Hempel '19 — Se n.d. pour les deux
  \fill[gray!13] (2,2) rectangle (3,3);
  \draw[black!15] (2,2) rectangle (3,3);
  \node[font=\tiny\bfseries] at (2.5, 2.72) {2 études};
  \node[font=\tiny, black!65] at (2.5, 2.50) {van~Hertem~'14};
  \node[font=\tiny, black!65] at (2.5, 2.28) {Hempel~'19};

  %% Multi-fermes : Thorup + Beer + Lavrova — Se 77–86 %
  \fill[rulegreen!22] (3,2) rectangle (4,3);
  \draw[black!18] (3,2) rectangle (4,3);
  \node[font=\tiny\bfseries] at (3.5, 2.72) {3 études};
  \node[font=\tiny] at (3.5, 2.50) {Se\,77–86\,\%};
  \node[font=\tiny, black!65] at (3.5, 2.26) {Thorup~\emph{et~al.}};

  %% Revue : González '23 — n.d.
  \fill[gray!13] (4,2) rectangle (5,3);
  \draw[black!15] (4,2) rectangle (5,3);
  \node[font=\tiny\bfseries] at (4.5, 2.65) {1 étude};
  \node[font=\tiny, black!65] at (4.5, 2.37) {González~'23};

  %% Synthétique : vide
  \fill[white] (5,2) rectangle (6,3);
  \draw[gray!20, dashed, thin] (5,2) rectangle (6,3);

  %% ==========================================================
  %% ROW 1 — APPRENTISSAGE PROFOND
  %% ==========================================================

  %% Conceptuel, 1 ferme, Terrain : vides
  \foreach \c in {0,1,2} {
    \fill[white] (\c,1) rectangle (\c+1,2);
    \draw[gray!20, dashed, thin] (\c,1) rectangle (\c+1,2);
  }

  %% Multi-fermes : Balasso '23 — Se = 85 %
  \fill[rulegreen!22] (3,1) rectangle (4,2);
  \draw[black!18] (3,1) rectangle (4,2);
  \node[font=\tiny\bfseries] at (3.5, 1.68) {1 étude};
  \node[font=\tiny] at (3.5, 1.45) {Se\,=\,85\,\%};
  \node[font=\tiny, black!65] at (3.5, 1.22) {Balasso~'23};

  %% Revue : Qiao '21 — n.d.
  \fill[gray!13] (4,1) rectangle (5,2);
  \draw[black!15] (4,1) rectangle (5,2);
  \node[font=\tiny\bfseries] at (4.5, 1.65) {1 étude};
  \node[font=\tiny, black!65] at (4.5, 1.37) {Qiao~'21};

  %% Synthétique : vide
  \fill[white] (5,1) rectangle (6,2);
  \draw[gray!20, dashed, thin] (5,1) rectangle (6,2);

  %% ==========================================================
  %% ROW 0 — NON SUPERVISÉ
  %% ==========================================================

  %% Conceptuel : lacune
  \fill[rawred!9] (0,0) rectangle (1,1);
  \draw[gray!25, dashed, thin] (0,0) rectangle (1,1);
  \node[black!58, font=\tiny\itshape] at (0.5, 0.5) {lacune};

  %% 1 ferme : Michelena '25 — n.d.
  \fill[gray!13] (1,0) rectangle (2,1);
  \draw[black!15] (1,0) rectangle (2,1);
  \node[font=\tiny\bfseries] at (1.5, 0.68) {1 étude};
  \node[font=\tiny, black!65] at (1.5, 0.38) {Michelena~'25};

  %% Terrain, Multi-fermes, Revue : lacunes
  \foreach \c in {2,3,4} {
    \fill[rawred!9] (\c,0) rectangle (\c+1,1);
    \draw[gray!25, dashed, thin] (\c,0) rectangle (\c+1,1);
    \node[black!58, font=\tiny\itshape] at (\c+0.5, 0.5) {lacune};
  }

  %% Synthétique : CE MÉMOIRE  (fond uniquement ici)
  \fill[rulegreen!48] (5,0) rectangle (6,1);
  \node[rulegreen!22!black, font=\tiny\bfseries, align=center]
      at (5.5, 0.72) {\textbf{Ce mémoire}};
  \node[rulegreen!22!black, font=\tiny, align=center]
      at (5.5, 0.50) {Se\;90–100\,\%};
  \node[rulegreen!22!black, font=\tiny, align=center]
      at (5.5, 0.27) {non supervisé};

  %% ===== CADRE EXTÉRIEUR + SÉPARATEURS =====
  \draw[black!35] (0,0) rectangle (6,4);
  \foreach \x in {1,2,3,4,5} \draw[black!18, thin] (\x,0) -- (\x,4);
  \foreach \y in {1,2,3} \draw[black!18, thin] (0,\y) -- (6,\y);

  %% Bordure CE MÉMOIRE redessinée EN DERNIER pour couvrir le cadre global
  \draw[rulegreen!60!black, line width=1.6pt] (5,0) rectangle (6,1);

  %% ===== LÉGENDE =====
  \fill[rulegreen!48] (0,-0.42) rectangle (0.52,-0.12);
  \draw[rulegreen!55!black, line width=0.9pt] (0,-0.42) rectangle (0.52,-0.12);
  \node[font=\tiny, anchor=west] at (0.58,-0.27) {Se\,$\geq$\,90\,\%};

  \fill[rulegreen!20] (1.55,-0.42) rectangle (2.07,-0.12);
  \draw[black!18] (1.55,-0.42) rectangle (2.07,-0.12);
  \node[font=\tiny, anchor=west] at (2.13,-0.27) {Se\,77–89\,\%};

  \fill[loforange!18] (3.15,-0.42) rectangle (3.67,-0.12);
  \draw[black!18] (3.15,-0.42) rectangle (3.67,-0.12);
  \node[font=\tiny, anchor=west] at (3.73,-0.27) {Se\,$<$\,77\,\%\;ou\;forte dispers.};

  \fill[gray!13] (0,-0.80) rectangle (0.52,-0.50);
  \draw[black!15] (0,-0.80) rectangle (0.52,-0.50);
  \node[font=\tiny, anchor=west] at (0.58,-0.65) {Se non disponible};

  \fill[rawred!9] (1.55,-0.80) rectangle (2.07,-0.50);
  \draw[gray!25, dashed, thin] (1.55,-0.80) rectangle (2.07,-0.50);
  \node[font=\tiny, anchor=west] at (2.13,-0.65) {Lacune (aucun travail publié)};

  \node[font=\tiny, black!58, anchor=west] at (3.8,-0.65)
      {$^{*}$ Se\,\% disponible pour une partie seulement des études};

\end{tikzpicture}%
}% fin resizebox
\caption{Matrice de couverture des travaux représentatifs en détection de boiterie bovine.}
    \label{fig:evidence_map}
\end{figure}

\FloatBarrier

% =====================================================================
\section{Synthèse critique}
% =====================================================================

La littérature et les guides pratiques indiquent trois exigences pour une solution déployable
de détection de boiterie:
\begin{enumerate}
  \item \textit{Sensibilité opérationnelle}: détecter les cas plausiblement détectables,
  en particulier les boiteries modérées à sévères qui nécessitent une intervention;
  \item \textit{Maîtrise des fausses alertes}: limiter les faux positifs et la fatigue
  d'alerte, qui dégradent rapidement la confiance des utilisateurs dans le système
  \citep{rutten2013invited};
  \item \textit{Interprétabilité}: fournir des raisons explicites et actionnables pour
  chaque alerte, afin de faciliter la prise de décision terrain.
\end{enumerate}

En pratique, les méthodes purement pilotées par les données peuvent manquer d'explicabilité pour
l'opérationnel, tandis que les approches purement à règles peuvent manquer de flexibilité
face à la variabilité inter-animaux. Un compromis robuste consiste à coupler détection
d'anomalies et contraintes métier explicites. Ce couplage cumule la
capacité de généralisation des algorithmes non supervisés et des gardes-fous
interprétables. La littérature récente confirme cette orientation: les systèmes
d'élevage de précision les plus prometteurs combinent des signaux hétérogènes avec des
logiques de décision interprétables, plutôt que de s'appuyer sur un modèle unique en boîte
noire \citep{palma2025scoping,michelena2025plf_review}.

L'hétérogénéité des protocoles évoquée à la section précédente a une conséquence directe
sur l'évaluation des systèmes: de nombreux travaux comparent des capteurs, des fermes, des
labels et des définitions de cas différents, ce qui complique l'appréciation d'une chaîne
déployable de bout en bout. Dans ce contexte, il existe une place claire pour des
contributions moins théoriques mais plus systémiques: un pipeline reproductible, une
validation explicite sur plusieurs niveaux et une traçabilité des artefacts expérimentaux.

La revue de littérature conduit donc à un enseignement central: la difficulté n'est pas
seulement de proposer un détecteur performant, mais de construire une chaîne de décision
cohérente, reproductible et défendable du point de vue métier. Une telle exigence justifie
la structure du présent mémoire, qui traite l'algorithme, les règles métier, l'interface et
la validation comme un ensemble cohérent plutôt que comme des blocs séparés.

% =====================================================================
\section{Positionnement du mémoire}
% =====================================================================

Le présent mémoire s'inscrit dans le compromis identifié ci-dessus, à savoir le couplage
entre détection d'anomalies non supervisée et règles métier interprétables. Sa contribution,
orientée ingénierie scientifique, vise à combler le manque identifié dans la littérature en
proposant un système validé de bout en bout, assorti d'une traçabilité expérimentale
complète. Concrètement, le mémoire livre:
\begin{enumerate}
  \item un pipeline complet depuis les données capteurs jusqu'aux notifications
  interprétables, combinant détection d'anomalies par Isolation Forest et règles métier
  explicites;
  \item une validation à deux niveaux, à savoir une validation manuelle par injections
  contrôlées, puis une validation robuste multi-configurations sur 9\,504~exécutions couvrant
  336~configurations;
  \item une étude d'ablation comparant quatre variantes du pipeline, incluant un détecteur
  alternatif (LOF);
  \item une traçabilité reproductible assurée par tests automatisés, comparaisons
  inter-versions et vérification d'intégrité cryptographique.
\end{enumerate}

La portée reste clairement délimitée: la sortie est une \textit{détection de boiterie
probable} et non un diagnostic clinique vétérinaire. Cette distinction est fondamentale et
cohérente avec les recommandations de la littérature \citep{alsaaod2019automatic}.

Le chapitre suivant détaille donc les données utilisées, les transformations appliquées aux
signaux capteurs, le choix des \textit{features} et le protocole expérimental retenu. L'objectif n'est
pas seulement de décrire une implémentation, mais de montrer comment la revue précédente se
traduit en décisions méthodologiques explicites.
